\documentclass[10pt, letterpaper]{article}

% Packages and styling copied from the original RenderCV template.  These
% settings provide the fonts, colors and custom environments used to
% structure entries throughout the document.
\usepackage[
    ignoreheadfoot,
    top=2 cm,
    bottom=2 cm,
    left=2 cm,
    right=2 cm,
    footskip=1.0 cm
]{geometry}
\usepackage{titlesec}
\usepackage{tabularx}
\usepackage{array}
\usepackage[dvipsnames]{xcolor}
% Primary colour: black for a clean professional look
\definecolor{primaryColor}{RGB}{0, 0, 0}
\usepackage{enumitem}
\usepackage{fontawesome5}
\usepackage{amsmath}
\usepackage[
    pdftitle={Zhaoxiang (Aaron) Feng's CV},
    pdfauthor={Zhaoxiang Feng},
    pdfcreator={LaTeX with RenderCV},
    colorlinks=true,
    urlcolor=primaryColor
]{hyperref}
\usepackage[pscoord]{eso-pic}
\usepackage{calc}
\usepackage{bookmark}
\usepackage{lastpage}
\usepackage{changepage}
\usepackage{paracol}
\usepackage{ifthen}
\usepackage{needspace}
\usepackage{iftex}

% Ensure machine readability
\ifPDFTeX
    \input{glyphtounicode}
    \pdfgentounicode=1
    \usepackage[T1]{fontenc}
    \usepackage[utf8]{inputenc}
    \usepackage{lmodern}
\fi
\usepackage{charter}

% Page and section formatting
\raggedright
\AtBeginEnvironment{adjustwidth}{\partopsep0pt}
\pagestyle{empty}
\setcounter{secnumdepth}{0}
\setlength{\parindent}{0pt}
\setlength{\topskip}{0pt}
\setlength{\columnsep}{0.15cm}
\pagenumbering{gobble}

\titleformat{\section}{\needspace{4\baselineskip}\bfseries\large}{}{0pt}{}[\vspace{1pt}\titlerule]
\titlespacing{\section}{-1pt}{0.3 cm}{0.2 cm}

\renewcommand\labelitemi{$\vcenter{\hbox{\small$\bullet$}}$}
\newenvironment{highlights}{\begin{itemize}[topsep=0.10 cm, parsep=0.10 cm, partopsep=0pt, itemsep=0pt, leftmargin=0 cm + 10pt]}{\end{itemize}}
\newenvironment{highlightsforbulletentries}{\begin{itemize}[topsep=0.10 cm, parsep=0.10 cm, partopsep=0pt, itemsep=0pt, leftmargin=10pt]}{\end{itemize}}
\newenvironment{onecolentry}{\begin{adjustwidth}{0 cm + 0.00001 cm}{0 cm + 0.00001 cm}}{\end{adjustwidth}}
\newenvironment{twocolentry}[2][]{\onecolentry \def\secondColumn{#2} \setcolumnwidth{\fill, 4.5 cm} \begin{paracol}{2}}{\switchcolumn \raggedleft \secondColumn \end{paracol} \endonecolentry}

\newenvironment{header}{\setlength{\topsep}{0pt}\par\kern\topsep\centering\linespread{1.5}}{\par\kern\topsep}

% Last updated text in the top right corner
\newcommand{\placelastupdatedtext}{%
  \AddToShipoutPictureFG*{%
    \put(\LenToUnit{\paperwidth-2 cm-0 cm+0.05cm},\LenToUnit{\paperheight-1.0 cm}){\vtop{\null\makebox[0pt][c]{\small\color{gray}\textit{Last updated in October 2025}}}}%
  }%
}

% Save original href for custom link formatting
\let\hrefWithoutArrow\href
% Redefine href to remove arrow icons
\renewcommand{\href}[2]{\hrefWithoutArrow{#1}{#2}}

\begin{document}
    % Draw the last updated text on each page
    \placelastupdatedtext

    % Header with personal information
    \begin{header}
        \textbf{\fontsize{24 pt}{24 pt}\selectfont Zhaoxiang (Aaron) Feng}

        \vspace{5 pt}

        \normalsize
        \mbox{{\color{black}\footnotesize\faMapMarker*}\hspace*{0.13cm}Del Mar, CA} \kern 5 pt \AND
        \mbox{\hrefWithoutArrow{mailto:zhf004@ucsd.edu}{\color{black}{\footnotesize\faEnvelope[regular]}\hspace*{0.13cm}zhf004@ucsd.edu}} \kern 5 pt \AND
        \mbox{\hrefWithoutArrow{tel:+13363377139}{\color{black}{\footnotesize\faPhone*}\hspace*{0.13cm}(336)\,337\,7139}} \kern 5 pt \AND
        \mbox{\hrefWithoutArrow{https://github.com/aaronzhfeng}{\color{black}{\footnotesize\faGithub}\hspace*{0.13cm}aaronzhfeng}} \kern 5 pt \AND
        \mbox{\hrefWithoutArrow{https://zhf004.github.io}{\color{black}{\footnotesize\faGlobe}\hspace*{0.13cm}zhf004.github.io}}
    \end{header}

    \vspace{5 pt - 0.3 cm}

    % Education
    \section{Education}

    \begin{twocolentry}{\textit{Sept.\,2022 -- Jun.\,2026 (expected)}}
        \textbf{University of California, San Diego (UCSD)}

        \textit{B.S. in Data Science \& B.S. in Probability \& Statistics}
    \end{twocolentry}

    \vspace{0.10 cm}
    \begin{onecolentry}
        \begin{highlights}
            \item GPA: 3.83 (Major GPAs: 3.94 in Data Science, 3.93 in Statistics)
            \item Honors: Provost Honors (multiple quarters)
            \item Graduate-level Coursework: DSC\,291 (Topics in Data Science), MATH\,287C (Advanced Time Series Analysis), DSC\,240 (Machine Learning), MATH\,282B (Applied Statistics\,II), MATH\,287A (Time Series Analysis), MATH\,278C (Seminar in Optimization), MATH\,282A (Applied Statistics\,I)
            \item Relevant Coursework: Machine Learning, Deep Learning, Natural Language Processing, Mathematical Statistics, Stochastic Processes, Data Management, Linear Algebra, Analysis
        \end{highlights}
    \end{onecolentry}

    % Publication
    \section{Publication}

    \begin{onecolentry}
        \textbf{ArcMemo: Abstract Reasoning Composition with Lifelong LLM Memory} \\
        \textit{Matthew Ho, Chen Si, Zhaoxiang Feng, Fangxu Yu, Yichi Yang, Zhijian Liu, Zhiting Hu, Lianhui Qin}. International Conference on Learning Representations (ICLR), 2026 (to appear). \\
        \href{ArcMemo_ICLR.pdf}{Paper} \\
        ArcMemo introduces a concept-level memory framework for large language models. By distilling reusable abstractions from reasoning traces, it enables continual learning and significantly outperforms strong no-memory baselines on the ARC-AGI benchmark.
    \end{onecolentry}

    % Research Experience
    \section{Research Experience}

    % Q-lab Research Assistant
    \begin{twocolentry}{\textit{Jan.\,2025 -- Present}}
        \textbf{Q-lab}\newline
        \textit{Research Assistant}
    \end{twocolentry}

    \vspace{0.10 cm}
    \begin{onecolentry}
        \begin{highlights}
            \item Developed \emph{ArcMemo}, a lifelong memory module for large language model reasoning. Built a pipeline to extract and abstract concepts from solution traces and integrate them into prompts, enabling continual improvement across tasks.
            \item Evaluated ArcMemo on the ARC‑AGI benchmark; achieved a substantial performance improvement over a strong no‑memory baseline and demonstrated that performance scales with additional inference compute.
            \item Designed dynamic feature detection routines and helper functions to generate transformation hints for compositional puzzles, and optimized asynchronous workflows for efficient model evaluation.
        \end{highlights}
    \end{onecolentry}

    \vspace{0.20 cm}

    % U.S. Immigration Policy Center
    \begin{twocolentry}{\textit{Sept.\,2023 -- Jun.\,2024}}
        \textbf{U.S. Immigration Policy Center}\newline
        \textit{Research Assistant}
    \end{twocolentry}

    \vspace{0.10 cm}
    \begin{onecolentry}
        \begin{highlights}
            \item Utilized Python and R to analyze historical naturalization data and forecast long‑term immigration trends, informing policy briefs for the 2024 election cycle.
            \item Developed reproducible data pipelines and visualization dashboards to communicate findings to non‑technical stakeholders.
        \end{highlights}
    \end{onecolentry}

    \vspace{0.20 cm}

    % Wang Lab
    \begin{twocolentry}{\textit{Jun.\,2025 -- Present}}
        \textbf{Wang Lab}\newline
        \textit{Student Researcher}
    \end{twocolentry}

    \vspace{0.10 cm}
    \begin{onecolentry}
        \begin{highlights}
            \item Working on computational chemistry projects exploring graph neural networks and transformer models; details forthcoming.
        \end{highlights}
    \end{onecolentry}

    \vspace{0.20 cm}

    % HDSI Student Researcher – Hao Zhang
    \begin{twocolentry}{\textit{Sept.\,2025 -- Present}}
        \textbf{Halıcıo\u{g}lu Data Science Institute, UC San Diego}\newline
        \textit{Student Researcher — Advisor: Prof. Hao Zhang}
    \end{twocolentry}

    \vspace{0.10 cm}
    \begin{onecolentry}
        \begin{highlights}
            \item Senior Capstone: “Open LLM Training, Inference, and Infrastructure” (ongoing). Building scalable pipelines for training open‑source large language models, benchmarking inference frameworks, and exploring alignment strategies.
        \end{highlights}
    \end{onecolentry}

    \vspace{0.20 cm}

    % HDSI Student Researcher – Tian Ma
    \begin{twocolentry}{\textit{Sept.\,2025 -- Present}}
        \textbf{Halıcıo\u{g}lu Data Science Institute, UC San Diego}\newline
        \textit{Student Researcher — Advisor: Prof. Tian Ma}
    \end{twocolentry}

    \vspace{0.10 cm}
    \begin{onecolentry}
        \begin{highlights}
            \item Senior Capstone: “Quantifying the credibility of large language model outputs” (ongoing). Investigating metrics and algorithms to measure factual consistency and reliability in LLM responses.
        \end{highlights}
    \end{onecolentry}

    % Teaching Experience
    \section{Teaching Experience}

    % Fall 2024
    \begin{twocolentry}{\textit{Fall 2024}}
        \textbf{Reader — MATH 180C: Introduction to Stochastic Processes II}\newline
        University of California, San Diego
    \end{twocolentry}

    \vspace{0.10 cm}
    \begin{onecolentry}
        \begin{highlights}
            \item Graded problem sets and exams, ensuring detailed feedback on Markov chains and stochastic calculus topics.
            \item Held weekly office hours to clarify course material and support student comprehension.
        \end{highlights}
    \end{onecolentry}

    \vspace{0.15 cm}

    % Winter 2025
    \begin{twocolentry}{\textit{Winter 2025}}
        \textbf{Reader — MATH 185: Introduction to Computational Statistics}\newline
        University of California, San Diego
    \end{twocolentry}

    \vspace{0.10 cm}
    \begin{onecolentry}
        \begin{highlights}
            \item Evaluated programming assignments and exams covering statistical simulation, Monte Carlo methods, and estimation techniques.
            \item Collaborated with instructors to provide clarifications on computational projects.
        \end{highlights}
    \end{onecolentry}

    \vspace{0.15 cm}

    % Spring 2025 MATH 180C
    \begin{twocolentry}{\textit{Spring 2025}}
        \textbf{Reader — MATH 180C: Introduction to Stochastic Processes II}\newline
        University of California, San Diego
    \end{twocolentry}

    \vspace{0.10 cm}
    \begin{onecolentry}
        \begin{highlights}
            \item Provided detailed feedback on assignments focusing on Poisson processes and discrete‑time martingales.
            \item Assisted with exam proctoring and maintained grade records.
        \end{highlights}
    \end{onecolentry}

    \vspace{0.15 cm}

    % Spring 2025 MATH 181A
    \begin{twocolentry}{\textit{Spring 2025}}
        \textbf{Reader — MATH 181A: Introduction to Mathematical Statistics I}\newline
        University of California, San Diego
    \end{twocolentry}

    \vspace{0.10 cm}
    \begin{onecolentry}
        \begin{highlights}
            \item Graded problem sets and exams on estimation, hypothesis testing, and confidence intervals.
            \item Supported the instructor by addressing student questions via discussion forums.
        \end{highlights}
    \end{onecolentry}

    \vspace{0.15 cm}

    % Summer Session I 2025
    \begin{twocolentry}{\textit{Summer Session I 2025}}
        \textbf{Tutor — DSC 40A: Theoretical Foundations of Data Science I}\newline
        University of California, San Diego
    \end{twocolentry}

    \vspace{0.10 cm}
    \begin{onecolentry}
        \begin{highlights}
            \item Conducted tutoring sessions covering set theory, probability, and algorithmic thinking for early‑career data science students.
            \item Led review sessions and prepared supplementary materials to aid student learning.
        \end{highlights}
    \end{onecolentry}

    \vspace{0.15 cm}

    % Summer Session II 2025
    \begin{twocolentry}{\textit{Summer Session II 2025}}
        \textbf{Reader — MATH 181A: Introduction to Mathematical Statistics I}\newline
        University of California, San Diego
    \end{twocolentry}

    \vspace{0.10 cm}
    \begin{onecolentry}
        \begin{highlights}
            \item Assisted with grading and feedback for advanced statistical inference coursework.
        \end{highlights}
    \end{onecolentry}

    \vspace{0.15 cm}

    % Fall 2025
    \begin{twocolentry}{\textit{Fall 2025}}
        \textbf{Reader — MATH 180A: Introduction to Probability (Sections B00 \& C00)}\newline
        University of California, San Diego
    \end{twocolentry}

    \vspace{0.10 cm}
    \begin{onecolentry}
        \begin{highlights}
            \item Managed grading for two lecture sections, covering probability spaces, random variables, and limit theorems.
            \item Provided consultation during office hours on problem‑solving strategies.
        \end{highlights}
    \end{onecolentry}

    % Selected Projects
    \section{Selected Projects}

    % 3D-Enhanced Graph Reaction Prediction
    \begin{twocolentry}{\textit{\href{https://aaronzhfeng.github.io/projects/3d-enhanced-graph-reaction}{3D‑Enhanced Graph Reaction Prediction}}}
        \textbf{3D‑Enhanced Graph Reaction Prediction}
    \end{twocolentry}

    \vspace{0.10 cm}
    \begin{onecolentry}
        \begin{highlights}
            \item Designed a graph neural network for reaction outcome prediction combining 2D molecular graph encoders with 3D geometry‑aware message passing via a contrastive InfoNCE objective.
            \item Pre‑trained 2D and 3D encoders on the QM9 dataset to infuse spatial knowledge and fine‑tuned on the Open Reaction Database (ORD); observed that incorporating 3D features improves representations for stereochemically constrained reactions.
            \item Integrated the graph encoder with an LSTM decoder to generate product SMILES, demonstrating improved accuracy on reactions with significant geometric constraints.
        \end{highlights}
    \end{onecolentry}

    \vspace{0.2 cm}

    % ChemTransformer
    \begin{twocolentry}{\textit{\href{https://aaronzhfeng.github.io/projects/chemtransformer}{ChemTransformer}}}
        \textbf{ChemTransformer}
    \end{twocolentry}

    \vspace{0.10 cm}
    \begin{onecolentry}
        \begin{highlights}
            \item Built a transformer‑based sequence‑to‑sequence model for chemical reaction prediction, leveraging SMILES and SELFIES tokenization and novel compound‑based relative positional encodings.
            \item Explored multiple model configurations on the Open Reaction Dataset (ORD) and found that traditional token‑level positional encoding outperformed compound‑level variants, while SELFIES improved output validity.
            \item Achieved competitive sequence‑level accuracy despite training challenges, highlighting the potential of transformer architectures for large‑scale reaction prediction.
        \end{highlights}
    \end{onecolentry}

    % Honors & Awards
    \section{Honors \& Awards}

    \begin{onecolentry}
        \begin{highlights}
            \item Provost Honors (Fall 2022 – Spring 2025)
            \item Conference Travel Funding — Halıcıo\u{g}lu Data Science Institute (Jul. 2025)
            \item Conference Travel Support — Sixth College at UC San Diego (Jul. 2025)
        \end{highlights}
    \end{onecolentry}

    % Technical Skills
    \section{Technical Skills}

    \begin{onecolentry}
        \textbf{Languages/Frameworks:} Python, Java, R, SQL, JavaScript, HTML \\
        \textbf{ML \& Data Tools:} PyTorch, scikit‑learn, Pandas, Dask, Spark, AWS \\
        \textbf{NLP/LLM:} Transformers, LLM fine‑tuning, retrieval‑augmented generation, agent building, GraphRAG, OpenAI/Anthropic APIs (sync/async requests) \\
        \textbf{Systems/Tooling:} Distributed training, Weights \& Biases, Docker, Linux, SSH
    \end{onecolentry}

\end{document}